\section{Objetivos y alcance}
\label{sec:objetivos}

\textbf{Objetivo general:} diseñar, simular y preparar para hardware SDR un enlace digital BPSK de telemetría industrial, con un receptor propio capaz de recuperar la información bajo un modelo de ruido AWGN más ruido impulsivo, y evaluar su desempeño mediante curvas BER frente a $E_b/N_0$.
    
\textbf{Objetivos específicos:}
\begin{enumerate}
    \item Relacionar los requerimientos de la aplicación con $R_s$, $R_b$, SPS, la tasa de muestreo y el filtro RRC.
    \item Diseñar el receptor (filtro acoplado, observación del ojo, muestreo y decisión) y validarlo en condición ideal y con el modelo de canal del proyecto.
    \item Cuantificar el efecto del ruido impulsivo en tiempo, frecuencia, constelación y tasa de error de bit.
    \item Configurar la interfaz con HackRF One y USRP sin alterar los parámetros de diseño y preparar la medición BER experimental.
\end{enumerate}

\textbf{Alcance:} el canal se modela en banda base con AWGN de amplitud 0.10 y un tren periódico de impulsos reales de amplitud $A_{imp}\in\{0,\,0.3,\,0.6,\,1.0\}$, un impulso cada 256 muestras; para el análisis de robustez se extiende a $A_{imp}=5$. En simulación se asume temporización ideal; en hardware se incorporan recuperación de frecuencia, reloj y fase. Las mediciones con el radio físico se realizan en el laboratorio con los flujogramas descritos.
